\section{Conclusion}
\label{sec:conclusion}

This report formulated operating-theatre planning as a stochastic program solved
through sample average approximation, progressing from operational scheduling
(Steps~1--3: timing, sequencing, and full assignment) through its statistical
validation (Step~4: SAA convergence, EVPI, and VSS) to a tactical three-stage
blueprint that fixes a per-session case-mix template before the waiting list is
known (Step~5).

At the \emph{operational} level
(Section~\ref{subsec:operational_insights}), waiting time dominates cost, so good
schedules build in early slack and trade a little idle time to prevent overrun
cascades; specialty mismatches arise only under a capacity shortfall, so blocks
should be kept specialty-pure; and the positive VSS with a large EVPI shows that
explicitly modelling duration uncertainty---rather than planning against mean
durations---drives the gains, most of all at the assignment level. At the
\emph{tactical} level (Section~\ref{sec:bp_results}), the blueprint is robust
across seeds: dedicate one session as an anchor for the long, high-variance
procedures, densify the others with short, predictable cases, provision at least
one long-case slot, and spread medium cases uniformly. Imposing this template
proved operationally close to free---cost differences sit within solver
gaps---provided the quotas carry a small safety margin, since worst-case-tight
provisioning can become infeasible when realised demand exceeds the sampled worst
case.

The two levels are complementary: the blueprint structures \emph{which} cases
meet \emph{which} session, while the operational model absorbs the residual
duration uncertainty within each session. The principal barrier to sharper
conclusions is the optimality gap left by the Big-$M$ formulation; closing it
(Section~\ref{sec:discussion}) is the most valuable next step.
